\documentclass[11pt]{article}
\usepackage{graphicx} % Required for inserting images
\usepackage{hyperref} % Required for hyperlinks

% Beautiful theorems, definitions etc.
\usepackage[most]{tcolorbox}
\usepackage{amsmath, amsthm}

% This command forces dotted lines for sections in the article class
\usepackage{tocloft}
\renewcommand{\cftsecleader}{\cftdotfill{\cftdotsep}}

\title{Linear Algebra and Optimization}
\author{Xavier Henry}
\date{September 2026}

% --- 1. Page Setup & Geometry ---
\usepackage[utf8]{inputenc}
\usepackage[margin=1in]{geometry} % Fixes the massive default LaTeX margins

% --- 2. Math & Science Essentials ---
\usepackage{amsmath, amssymb, amsthm} % Essential for formulas, symbols, and proofs
\usepackage{bm} % For bold math symbols
\usepackage{enumitem} % for lists


% --- 3. Beautiful Callout Boxes (Perfect for Notes!) ---
% 3.1. Define a standard theorem counter that resets with the section
\newtheorem{thm}{Theorem}[section]
\newtheorem{defnt}{Definition}[section]
\newtheorem{exmpl}{Example}[section]
\newtheorem{corol}{Corollary}[section]

% 3.2. Create the color boxes and link them to those counters using "use counter"
\newtcolorbox[use counter=defnt]{definitionbox}[1]{
    colback=blue!5!white, 
    colframe=blue!75!black, 
    fonttitle=\bfseries, 
    title=Definition \thedefnt: #1, % Automatically prints "Definition X.Y: Your Title"
    arc=4mm
}

\newtcolorbox[use counter=thm]{theorembox}[1]{
    colback=purple!5!white,    % 5% purple background (soft pastel)
    colframe=purple!75!black,  % Deep purple border and title bar 
    fonttitle=\bfseries, 
    % This line checks if you gave a title. If not, it skips the colon
    title={Theorem \thethm\ifstrempty{#1}{}{ : #1}},
    arc=4mm
}

\newtcolorbox[use counter=exmpl]{examplebox}[1]{
    colback=green!3!white,    % 5% purple background (soft pastel)
    colframe=green!80!black,  % Deep purple border and title bar 
    fonttitle=\bfseries, 
    % This line checks if you gave a title. If not, it skips the colon
    title={Example \theexmpl\ifstrempty{#1}{}{ : #1}},
    arc=4mm
}

\newtcolorbox[use counter=corol]{corollarybox}[1]{
    colback=yellow!5!white,    % 5% purple background (soft pastel)
    colframe=yellow!75!black,  % Deep purple border and title bar 
    fonttitle=\bfseries, 
    % This line checks if you gave a title. If not, it skips the colon
    title={Corollary \thecorol\ifstrempty{#1}{}{ : #1}},
    arc=4mm
}

\begin{document}

\maketitle

\tableofcontents \newpage

\section{Vectors, Matrices and Linear Equations}

\begin{definitionbox}{Linear Function}
    A function $f$ that maps $n$ dimensional vectors to $m$ dimensional vectors is called \textbf{linear} if it satisfies the properties:
    \begin{enumerate}
        \item For all vectors $x$ and all scalars $\alpha$, $f(\alpha x) = \alpha f(x)$.
        \item For all vectors $x$ and $y$, $f(x+y)=f(x)+f(y)$
    \end{enumerate}
\end{definitionbox}

For example, take $f$ to be the function that multiplies the entries of $x$ by 2. In this case, we can easily see $f$ is linear. The challenge arises from more exotic functions, such as swapping 2 coordinates of $x$, or adding one coordinate to all the others. It is important to note that \textbf{every linear function can be represented as a matrix}.

\begin{definitionbox}{Matrix-vector product}
    The matrix-vector product of a 2 x 2 matrix and a 2 dimensional vector is defined as
    \begin{align*}
        \begin{bmatrix}
            a & b \\
            c & d
        \end{bmatrix}
        \begin{bmatrix}
            x \\
            y
        \end{bmatrix}
        =
        \begin{bmatrix}
            ax+by \\
            cx+dy
        \end{bmatrix}
    \end{align*}
    We may also interpret this product as a linear combination of the columns of the matrix, which makes it easier to work with higher dimensions.
    \begin{align*}
        \begin{bmatrix}
            a & b & c \\
            d & e & f \\
            g & h & i
        \end{bmatrix}
        \begin{bmatrix}
            x \\
            y \\
            z
        \end{bmatrix}
        =
        x
        \begin{bmatrix}
            a \\
            d \\
            g
        \end{bmatrix}
        +
        y
        \begin{bmatrix}
            b \\
            e \\
            h
        \end{bmatrix}
        +
        z
        \begin{bmatrix}
            c \\
            f \\
            i
        \end{bmatrix}
    \end{align*}
    Moreover, we can interpret the result as a column matrix, where each entry is the dot product of the $n \times m$ matrix and our vector.
\end{definitionbox}

\begin{examplebox}{}
    What does the function that corresponds to multiplying a vector by the following matrix do to a point?
    \begin{align*}
        \begin{bmatrix}
            0 & 1 \\
            -1 & 0
        \end{bmatrix}
    \end{align*}
\end{examplebox}

\begin{proof}
    Take the product of our matrix and a generic 2 coordinate vector.
    \begin{align*}
        \begin{bmatrix}
            0 & 1 \\
            -1 & 0
        \end{bmatrix}
        \begin{bmatrix}
            x \\
            y
        \end{bmatrix}
        =
        \begin{bmatrix}
            y \\
            -x
        \end{bmatrix}
    \end{align*}
    If we graph this visually with any vector, we see that the function corresponds to rotating every vector in the $x-y$ plane clockwise by $\pi/2$.
\end{proof}

\begin{examplebox}{}
    What if we want to rotate every vector in the $x-y$ plane counterclockwise by $\pi/4$?
\end{examplebox}

\begin{proof}
    By far the easiest way to think of these (excluding just knowing it by heart which comes with practice) is to use the standard rotation matrix! The standard matrix for \textbf{counterclockwise} rotation is:

    \begin{align*}
        \begin{bmatrix}
            \cos\theta & -\sin\theta \\
            \sin\theta & \cos\theta
        \end{bmatrix}
    \end{align*}
    
    And now we can just plug in our desired angle. In this case, it's $\pi/4$:

    \begin{align*}
        \begin{bmatrix}
            \frac{\sqrt{2}}{2} & -\frac{\sqrt{2}}{2} \\
            \frac{\sqrt{2}}{2} & \frac{\sqrt{2}}{2}
        \end{bmatrix}
    \end{align*}
\end{proof}

\begin{theorembox}{Matrix-Vector Products are Linear Functions}
    We claim that any matrix-vector product is a linear function. Let's check this.
\end{theorembox}

\begin{proof}
    Take any scalar $\alpha$.
    \begin{align*}
        f\left(\alpha
        \begin{bmatrix}
            x \\ 
            y
        \end{bmatrix}
        \right)
        =
        \begin{bmatrix}
            a & b \\
            c & d
        \end{bmatrix}
        \begin{bmatrix}
            \alpha x \\
            \alpha y
        \end{bmatrix}
        =
        \begin{bmatrix}
            \alpha a x + \alpha b y \\
            \alpha c x + \alpha d y
        \end{bmatrix}
        =
        \alpha f\left(
        \begin{bmatrix}
            x \\
            y
        \end{bmatrix}
        \right)
    \end{align*}
    And take any two 2 vectors as shown below.
    \begin{align*}
        f\left(
            \begin{bmatrix}
                x \\
                y
            \end{bmatrix}
            +
            \begin{bmatrix}
                x' \\
                y'
            \end{bmatrix}
        \right)
        =
        \begin{bmatrix}
            a & b \\
            c & d
        \end{bmatrix}
        \begin{bmatrix}
            x+x' \\
            y+y'
        \end{bmatrix}
        =
        \begin{bmatrix}
            ax + by + ax' + by' \\
            cx + dy + cx' + dy'
        \end{bmatrix}
        =
        f\left(
            \begin{bmatrix}
                x \\
                y
            \end{bmatrix}
        \right)
        +
        f\left(
            \begin{bmatrix}
                x' \\
                y'
            \end{bmatrix}
        \right)
    \end{align*}
\end{proof}

\begin{definitionbox}{General Matrix-Vector Product}
    If $A$ is an $n \times m$ dimensional matrix and $x$ is an $m$ dimensional vector then $Ax$ is an $n$ dimensional vector and its ith coordinate is
    \begin{align*}
        (Ax)_i = \sum_{j=1}^{m} A_{i,j} x_j
    \end{align*}
\end{definitionbox}

\begin{theorembox}{Linear Functions are Matrix-Vector Products}
    Any linear function $f$ on vectors can be represented as a matrix-vector product. That is, there is some $A$ such that $f(x) = Ax$
\end{theorembox}


\begin{proof}
    Suppose a vector $x \in \mathbb{R}^m$ and linear function $f$.
   \begin{align*}
        \begin{bmatrix}
            x_1 \\
            x_2 \\
            \vdots \\
            x_m
        \end{bmatrix}
    \end{align*}
    We will construct a matrix $A \in \mathbb{R}^{m\times m}$ such that $f(x)=Ax$.
    We may write $x$ as a linear combination as follows
    \begin{align*}
        x_1
        \begin{bmatrix}
            1 \\
            0 \\
            \vdots \\
            0
        \end{bmatrix}
        +
        x_2
        \begin{bmatrix}
            0 \\
            1 \\
            \vdots \\
            0
        \end{bmatrix}
        +
        ...
        +
        x_m
        \begin{bmatrix}
            0 \\
            0 \\
            \vdots \\
            1
        \end{bmatrix}
    \end{align*}
    So we have
    \begin{align*}
        x =
        \sum_{i=1}^{m} x_i e_i
    \end{align*}
    where $e_i$ are the standard basis vectors in $\mathbb{R}^m$.
    \begin{align*}
        f(x) &=
        f\left(
            \sum_{i=1}^{m} x_i e_i
        \right) \\
        &=
        \sum_{i=1}^{m} x_i f(e_i)
    \end{align*}
    Let $A$ be denoted by
    \begin{align*}
        A &=
        \begin{bmatrix}
            f(e_1) & f(e_2) & ... & f(e_m)
        \end{bmatrix} \\
        &=
        \begin{bmatrix}
            a_{11} & ... & a_{1m} \\
            \vdots & \ddots & \vdots \\
            a_{m1} & ... & a_{mm}
        \end{bmatrix}
    \end{align*}
    By the definition of the matrix-vector product we have
    \begin{align*}
        Ax&=
        \begin{bmatrix}
            a_{11} & ... & a_{1m} \\
            \vdots & \ddots & \vdots \\
            a_{m1} & ... & a_{mm}
        \end{bmatrix}
        \begin{bmatrix}
            x_1 \\
            x_2 \\
            \vdots \\
            x_m
        \end{bmatrix} \\
        &=
        x_1
        \begin{bmatrix}
            a_{11} \\
            \vdots \\
            a_{m1}
        \end{bmatrix}
        +x_2
        \begin{bmatrix}
            a_{12} \\
            \vdots \\
            a_{m2}
        \end{bmatrix}
        + ...
        +x_m
        \begin{bmatrix}
            a_{1m} \\
            \vdots \\
            a_{mm}
        \end{bmatrix} \\
        &=
        \sum_{i=1}^{m} x_i f(e_i) \\
        Ax&=
        f(x)
    \end{align*}

\end{proof}
\newpage

\section{Systems of Linear Equations: Geometry and Solving}

\begin{definitionbox}{Inner Product}
    The inner product between two vectors $a$ and $b$ of dimension $n$ is
    \begin{align*}
        \langle a, b \rangle = \sum_{i=1}^{n} a_i b_i
    \end{align*}
    More commonly, this is called the dot product and can be written as $a \cdot b$. We can also write this as a matrix-vector product between a $1 \times n$ matrix and an $n$ dimensional vector. 
\end{definitionbox}

\begin{examplebox}{}
    Find the 3 by 3 matrix $R$ corresponding to a rotation of $60$ degrees around the y axis.
\end{examplebox}

\begin{examplebox}{}
    Recall that a degree 3 polynomial can be represented with a 4 dimensional vector. Find the matrix
    corresponding to a derivative in $\mathbb{R}^4 to \mathbb{R}^3$ and vice versa for an integral. and also an integral $\mathbb{R}^4 to \mathbb{R}^3$. 
\end{examplebox}

Suppose we start with two equations: $2x-3y=0$ and $x+y=5$. There are various ways to interpret this.

\subsection{Row Interpretation}
Each row gives a line.
\begin{align*}
    \begin{bmatrix}
        2 & -3
    \end{bmatrix}
    \begin{bmatrix}
        x \\
        y
    \end{bmatrix}
    =
    2x-3y=0 \\
    \begin{bmatrix}
        1 & 1
    \end{bmatrix}
    \begin{bmatrix}
        x \\
        y
    \end{bmatrix}
    =
    x + y = 5
\end{align*}
In this case, we have a solution, since these lines have unequal gradients. This is the expected scenario when the number of variables is equal to the number of constraints. Alternatively, if we have more variables than constraints, we would expect infinitely many solutions. If we had more constraints than variables, we will \textbf{likely?} have no solution. \\
We also have infinitely many solutions if the lines are the same, or the equations always hold true. For example, $0x+0y=0$ holds $\forall x, y$

\subsection{Column Interpretation}
Are there $x,y$ such that
\begin{align*}
    x
    \begin{bmatrix}
        2 \\
        1
    \end{bmatrix}
    + y
    \begin{bmatrix}
        -3 \\
        1
    \end{bmatrix}
    =
    \begin{bmatrix}
        0 \\
        5
    \end{bmatrix}
\end{align*}
It's helpful to graph this on the cartesian plane and see how operations on the vectors affect our visual outcomes.

\begin{theorembox}{Solving 2x2 Systems}
    The system of equations
    \begin{align*}
        \begin{cases}
            ax+by=\alpha \\
            cx+dy=\beta
        \end{cases}
    \end{align*}
    has a unique solution whenever $ad-bc \neq 0$.
\end{theorembox}
Below are three ways to prove this. You can probably tell the one I struggled with...

\subsection{Geometric Proof}
\begin{proof}
    The slopes of the lines are $\frac{-a}{b}$ and $\frac{-c}{d}$. The system has exactly one solution when their slopes are unequal, i.e. $ad-bc \neq 0$.
\end{proof}

\subsection{Symbolic Proof}
\begin{proof}
    Multiply the first equation by $d$ and the second by $b$ then the second by $a$ and the first by $c$.
    \begin{align*}
        adx+bdy-bcx-bdy&=d\alpha - b\beta \\
        (ad-bc)x=d\alpha -b\beta \\
    \end{align*}
    and also
    \begin{align*}
        (ad-bc)y=a\beta -c\alpha
    \end{align*}
    And $x,y \in \mathbb{R}$ iff $ad-bc \neq 0$
\end{proof}

\subsection{Algrebraic Proof}
Before we get to this proof, it's important to establish a few definitions.

\begin{definitionbox}{Onto}
    A function is onto if every element in the codomain is used as an output at least once.
\end{definitionbox}

\begin{theorembox}{Shoelace Theorem}
    Given a polygon with vertices $(x_1,y_1), (x_2,y_2),..., (x_n,y_n)$, listed in clockwise order, the area of such polygon is
    \begin{align*}
        A = \frac{1}{2} |(x_1y_2+x_2y_3+...+x_ny_1)-(y_1x_2+...+x_ny_1)|
    \end{align*}
\end{theorembox}

\begin{proof}
    Let
    \begin{align*}
        A =
        \begin{bmatrix}
            a & b \\
            c & d
        \end{bmatrix} \\
        B = 
        \begin{bmatrix}
            \alpha \\
            \beta
        \end{bmatrix}
    \end{align*}
    We view $A$ as a linear transformation to show that $A$ is both one-to-one and onto. Take a parallelogram with vertices $(0,0), (a,c), (b,d), (a+b,c+d)$. We can verify by the shoelace theorem that the area of the parallelogram is $|ad-bc|$. Two of the edges of the parallelogram correspond to the two columns of $A$. Thus, every point inside the parallelogram is some linear combination of the first two columns. Note that we exclude the top and right sides for uniqueness. Since $B \in \mathbb{R}^2$, $B$ is a point within the parallelogram of form
    \begin{align*}
        \lambda
        \begin{bmatrix}
            a \\
            c
        \end{bmatrix}
        + \mu
        \begin{bmatrix}
            b \\
            d
        \end{bmatrix}
    \end{align*}
    where $0 \leq lambda, \mu \leq N$ for $N$ sufficiently large. We know every point lying within is uniquely expressed as a linear combination of our column matrices (one-to-one and onto). Thus, there is exactly one solution when $ad-bc \neq 0$.
\end{proof}

\subsection{More on the Geometry of Linear Equations}
\begin{definitionbox}{Affine Hyperplane}
    An affine hyperplane is the set of all solutions to a linear equation
    \begin{align*}
        \langle a,x \rangle = 
        a_1x_1 + a_2x_2 + ... + a_nx_n = b
    \end{align*}
\end{definitionbox}
Now the set of solutions to a system of linear equations $Ax=b$ is the intersection of affine hyperplanes.
\begin{theorembox}{Set of Solutions is the Intersection of Hyperplanes}
    Let
    \begin{align*}
        H_i = \{x \text{s.t.} A_{i,1}x_1 + A_{i,2}x_2 + ... + A_{i,n}x_n = b_i \}
    \end{align*}
    Where $H_i$ is the hyperplane correpsonding to the i'th row of $A$ and $b$. $x$ is a solution to $Ax=b$ iff 
    \begin{align*}
        x \in \bigcap_{i=1}^{m} H_i
    \end{align*}
\end{theorembox}

\begin{definitionbox}{Span of Vectors}
    The span of vectors $v_1,v_2,...,v_n$ is all the vectors that can be obtained as linear combinations of $v_1,v_2,...,v_n$.
\end{definitionbox}
Moreover, we can deduce:
\begin{theorembox}{}
    If $A$ is $m \times n$, the system of linear equations $Ax=b$ has a solution for every choice of $b$ iff the span of the columns of $A$ is all of $\mathbb{R}^m$.
\end{theorembox}
We can use this theorem to show some pretty interesting results. Consider the system of linear equations
\begin{examplebox}{}
    Does this system of linear equations have any solutions? If so how many?
    \begin{align*}
        \begin{bmatrix}
            2 & -3 & 0 \\
            1 & 1 & 2 \\
            3 & -2 & 2
        \end{bmatrix}
        \begin{bmatrix}
            x \\
            y \\
            z 
        \end{bmatrix}
        =
        \begin{bmatrix}
            1 \\
            1 \\
            3
        \end{bmatrix}
    \end{align*}
\end{examplebox}
\begin{proof}
    None of the rows of $A$ are scalar multiples of each other, but we can add the first two equations to get an important constraint
    \begin{align*}
        3x-2y+2z=2
    \end{align*}
    which actually contradicts our third equation, so there is no solution for this system. Because we can express the third equation as a linear combination of the first two columns, the span of the columbs of $A$ cannot be all of $\mathbb{R}^3$.
\end{proof}
\begin{definitionbox}{Linear Independence}
    A set of vectors $v_1,v_2,...,v_n$ is linearly independent if no $v_i$ can be expressed as a linear combination of the others. If $A$ is $n \times n$, the system of linear equations $Ax=b$ has a solution $\forall b$ iff the columns of $A$ are linearly independent.
\end{definitionbox}

\section{Gaussian Elimination}
\begin{definitionbox}{Augmented Matrix}
    The augmented matrix of a system of linear equations $Ax=b$ is of the form
    \begin{align*}
        [A|b]
    \end{align*}
\end{definitionbox}
\begin{definitionbox}{Pivot}
    The pivot in a row is the leftmost nonzero.
\end{definitionbox}

The key of Gaussian elimination is that we need 3 invertible operations:
\begin{itemize}
    \item scale a row by some nonzero number
    \item add a multiple of a row to another row
    \item swap rows
\end{itemize}
Our goal is to have our pivots left to right strictly. That is in \textbf{Row Echelon Form}.
\begin{definitionbox}{Singular Matrix}
    A matrix $A$ is singular if its row echelon form has a row full of zeroes. This means there is either no solution or an infinite number of solutions.
\end{definitionbox}
A variable with no pivot is free. That is, it can be set to anything. The number of dimensions of our solution space is equal to the number of free variables! This is a very cool and important fact to remember!!!!

\section{Products, Adjacency Matrices, Matrix Exponentiation}
\begin{examplebox}{}
    \begin{align}
        x_1+x_2=12 \\
        x_1=x_3+x_4 \\
        x_2+x_3=x_5 \\
        x_4+x_5=12
    \end{align}
\end{examplebox}
\begin{proof}
    We begin with the augmented matrix.
    \begin{align*}
        \begin{bmatrix}
            1 & 1 & 0 & 0 & 0 & 12 \\
            1 & 0 & -1 & -1 & 0 & 0\\
            0 & 1 & 1 & 0 & -1 & 0\\
            0 & 0 & 0 & 1 & 1 & 12
        \end{bmatrix}
    \end{align*}
    And after changing to row echelon form (REF) we see that the last row is singular (filled with zeroeos). The singularity of this matric leads to an infinite number of solutions (courtesy of our free variables). 
\end{proof}

\subsection{Matrix Multiplication}
$A(Bv)=Cv \implies C=AB$
Suppose we have an $m \times n$ matrix A and $n \times p$ matrix B. Then
\begin{align*}
    C_{i,k} &= \sum_{j=1}^{n} A_{i,j}B_{j,k} \\
            &= ith row A* kth col B \\
            &= C_k = AB_k
\end{align*}
The last line should be understood as the kth column of $C$.
\begin{theorembox}{Properties of Matrix Multiplication}
    We have the associative property such that $A(BC)=(AB)C$. We have the distributive property (which holds for the right and left) $A(B+C)=AB+AC$. It is important to note that it is not necessarily true that $AB = BA$. In fact, it is more often true that $AB \neq BA$.
\end{theorembox}

\begin{definitionbox}{Walk}
    A walk is a sequence of vertices such that every consecutive pair is connected by an edge, with repeats permitted. We can use an adjacency matrix to represent a directed (or undirected) graph of several nodes. Recall Markov Chains from $18.600$!
\end{definitionbox}

\begin{theorembox}{Number of n length walks}
    Suppose an adjacency matrix $A$. $A^n$ gives the number of $n$ length walks that exist from one vertex to another (or itself) for each vertex. $(A^n)_{i,j}$ is the $n$ length walk from $i$ to $j$.
\end{theorembox}
\begin{proof}
    TODO: write a proof for the preceding theorem
    We induct on $n$. We first prove the simplest case that $A^2$ gives the number of length $2$ walks between vertices for each vertex.
\end{proof}

\begin{definitionbox}{Transpose}
    The transpose of an $m \times n$ matrix $A$ denoted by $A^T$ is an $n \times m$ matrix with
    \begin{align*}
        (A^T)_{i,j}=A_{j,i}
    \end{align*}
    If $A=A^T$ then we say that $A$ is symmetric. We can also express the inner-product using the
    transpose of a vector
    \begin{align*}
        \langle a,b \rangle = a^Tb
    \end{align*}
\end{definitionbox}

\begin{definitionbox}{Outer Product}
    The outer product between two $n$ dimensional vectors $a$ and $b$ is an $n \times n$ matrix $ab^T$.
    This gives us another interpretation of matrix multiplication. If $A$ and $B$ are $m \times n$ and
    $p \times n$ matrices with columns $a_1,a_2,...,a_n$ and $b_1,b_2,...,b_n$ respectively then
    \begin{align*}
        AB^T = \sum_{i=1}^{n} a_ib_i^T
    \end{align*}
    It becomes easier to see when we think of it visually. Take $a_i$ to be an $m \times 1$ column
    vector. Then take $b_i^T$ to be a $1 \times p$ row vector. Then each multiplication of these vectors
    produces an $m \times p$ matrix. We then sum all these resulting matrices to get our final answer.
\end{definitionbox}

\subsection{Inverse of a Matrix}
\begin{definitionbox}{Right Inverse of a Matrix}
    Given a square matrix $A$, its right inverse is a matrix $A^-1$ with the property that
    \begin{align*}
        AA^(-1)=I
    \end{align*}
\end{definitionbox}
While matrix multiplication is in general not commutative, the left and right inverse of a matrix are
always the same
\begin{definitionbox}{Left Inverse of a Matrix}
A matrix $A^{-1}$ is a left inverse, i.e. $A^{-1}A=I$ if and only if it is the right inverse, i.e. 
    $AA^{-1}=I$.
\end{definitionbox}

Suppose we want to solve a linear system $Ax=b$. If we know the inverse $A^{-1}$, we can get a closed
form expression for the solution. Consider $x=A^{-1}b$.
\begin{align*}
    Ax = A(A^{-1}b) = (AA^{-1})b = Ib = b
\end{align*}

\section{Subspaces and Nullspaces}
\begin{definitionbox}{Span}
    The span of $v_1, v_2, ..., v_n$ where $v_i \in \mathbb{R}^m$ is all linear combinations in
    $\mathbb{R}^m$ that can be formed of all $n$ vectors.
\end{definitionbox}

\begin{definitionbox}{Vector Space}
    A vectors space is a set $V$ of objects closed under addition and scalar multiplication (closed
    under linear combinations). That is $\forall v,w \in V, v+w \in V$. $\forall v \in V, \alpha \in
    \mathbb{R}, \alpha v \in V$
\end{definitionbox}

\begin{definitionbox}{Subspace}
    A subspace $S \subseteq V$ is also closed ... (similar to above). A set of vectors
    ${v_1,..v_\lambda}$ generate $S$ if $\forall v_i \in S$ and span ${v_1,...v_\lambda} = S$
\end{definitionbox}

\begin{theorembox}{Subspces}
    For subspaces $S_1$ and $S_2$, $S_1 \bigcap S_2$ is a subspace. $S_1 \bigcup S_2$ is not
    necessarily a subspace.
\end{theorembox}

\begin{definitionbox}{Column Space}
    Given a matrix $A$, the column space of $A$ is the span of its columns. That is for $m \times n$
    matrix $A$ we have
    \begin{align*}
        \begin{bmatrix}
            a_1 & a_2 & ... & a_n
        \end{bmatrix} \\
    \end{align*}
    $C(A)=span({a_1,...a_n})={Ac | c \in \mathbb{R}^n}$
\end{definitionbox}

\begin{theorembox}{}
    $C(A), N(A)$ are vector spaces ($N(A)$ is the nullspace of A).
\end{theorembox}

\begin{theorembox}{}
    For $A_{m \times n}$, $C(A) \subseteq \mathbb{R}^m, N(A) \subseteq \mathbb{R}^n$ and vice versa
    $A^T$
\end{theorembox}

\begin{definitionbox}{Nullspace}
    The nullspace of $A$ is the set of all vectors $x$ that make the equation $Ax=0$ true.
\end{definitionbox}

\begin{theorembox}{}
    Consider $Ax=b$. \\
    If $b \notin C(A)$ then no solution.
    If $b \in C(A)$ then we have to consider two cases:
    \begin{itemize}
        \item $N(A)=\{0\} \implies$ 1 solution.
        \item $N(A)=\{0,x_1,...\} \implies$ infinite solutions.
    \end{itemize}
\end{theorembox}

\begin{theorembox}{Long ass Theorem}
    For $A_{n \times n}$, $A$ invertible $\iff$ $C(A)=\mathbb{R}^n \iff$ cols of $A$ are linear
    independent $\iff N(A)=\{0\} \iff Ax=b$ uniquely solvable $\forall b$.
\end{theorembox}

\begin{theorembox}{Important theorem}
    Gauss-Jordan elimination preserves nullspace at every step. It is important to note that
    Gauss-Jordan does NOT preserve the column space of $A$. \\
    For $B$ invertible and $A=BA'$ where
    $A'$ is some matrix, $N(A')=N(A)$. If $B$ not invertible $N(A') \leq N(A)$
\end{theorembox}

\begin{theorembox}{}
    For $T$ invertible, $N(A)=N(AT)$ \\
    $x \in N(TA) \iff TAx=0 \iff T^-1 T Ax = T^-1 0 \iff Ax=0$
\end{theorembox}

\begin{theorembox}{}
    For $A_{n \times n}, B_{n \times n}$, if $x \in N(A)$ and $x \in N(B), x \in N(A+B)$
\end{theorembox}

\begin{theorembox}{}
    $C(AB) \subseteq C(A)$
\end{theorembox}
\begin{proof}
    Exercise!   
\end{proof}

\end{document}
